\chapter{Einleitung}

Um den Einfall von Myonen als geladenene kosmische Strahlung auf die Erdoberfläche zu charakterisieren, wird in diesem Versuch ihre Winkelverteilung mit einem gasgefüllten Teilchendetektor gemessen. Dabei wird als Detektor die Prototyp-Driftkammer des B-GOOD-Experiments \cite[1]{515versuchsanleitung} mit zwei zueinander versetzten Schichten von hexagonalen Drahtzellen verwendet, welche mit einem Szintillationsdetektor als Triggersignal sowie einer Ausleseelektronik verbunden ist \cite[18-20]{hammanndiplom}.

Die Driftkammer wird zunächst mithilfe einer Probe von $\ce{^90Sr}$, welches $\beta$-Strahlung ($\beta$-Zerfall: $\ce{n -> p + e- + \bar{\nu}}$ \cite[4]{leotechniques}) emittiert \cite{nudatsr90}, kalibriert. Die so emittierten Elektronen sind gleich geladen wie die zu messenden Myonen und können aufgrund dann der großen Aktivität (verglichen mit der deutlichen geringeren Flussrate der Myonen) der \ce{^90Sr}-Quelle zur praktischen Detektorkalibration genutzt werden.
